% Esfera celeste — los dos números que ubican cualquier astro.
%
% Diagrama propio en TikZ: sin dependencias externas, sin licencias de terceros,
% vectorial y editable. Usa los colores del template (milcGradeDark toma el
% color de la línea de investigación, así que el diagrama se tiñe solo).
%
% Se incrusta vía \guiaDiagrama desde el bloque `recursos:` del YAML.
\begin{tikzpicture}[scale=0.92,>=stealth,font=\sffamily]

  % ── Plano del horizonte (elipse en perspectiva) ──────────────────────
  \draw[thick,milcNegro!40] (0,0) ellipse (3.7 and 1.2);

  % ── Cúpula celeste ───────────────────────────────────────────────────
  \draw[thick,milcNegro!40] (-3.7,0) arc (180:0:3.7 and 3.2);

  % ── Observador ───────────────────────────────────────────────────────
  \fill[milcNegro] (0,0) circle (2.2pt);
  \node[font=\scriptsize\bfseries,below left=-1pt and 0pt] at (0,0) {Tú};

  % ── Cénit ────────────────────────────────────────────────────────────
  \fill[milcNegro!70] (0,3.2) circle (1.6pt);
  \node[font=\scriptsize\bfseries,above] at (0,3.2) {Cénit · altura $90^\circ$};
  \draw[dotted,milcNegro!35] (0,0) -- (0,3.2);

  % ── Puntos cardinales (el azimut se cuenta desde el N hacia el E) ────
  \node[font=\scriptsize\bfseries,above] at (0,1.2) {N $0^\circ$};
  \node[font=\scriptsize\bfseries,right] at (3.7,0) {E $90^\circ$};
  \node[font=\scriptsize\bfseries,below] at (0,-1.2) {S $180^\circ$};
  \node[font=\scriptsize\bfseries,left] at (-3.7,0) {O $270^\circ$};

  % ── Una estrella de ejemplo ──────────────────────────────────────────
  \coordinate (Pie) at (2.15,-0.62);   % su proyección en el horizonte
  \coordinate (Est) at (2.15,1.78);    % la estrella

  \draw[milcNegro!30] (0,0) -- (Est);              % línea de visión
  \draw[dashed,milcNegro!45] (Pie) -- (Est);       % vertical al horizonte

  % Azimut: arco sobre el horizonte, del N al pie de la estrella.
  % Etiqueta pegada al arco (abajo-izquierda), lejos de la de altura.
  \draw[->,line width=1pt,milcGradeDark]
    (0,1.2) .. controls (1.7,1.02) and (2.35,0.18) .. (Pie);
  \node[font=\scriptsize\bfseries,milcGradeDark,align=center] at (1.28,0.02)
    {azimut $\approx 115^\circ$};

  % Altura: del horizonte a la estrella.
  % Etiqueta a la derecha de la flecha vertical, sin chocar con el azimut.
  \draw[->,line width=1pt,milcTurquesa] (Pie) -- (Est);
  \node[font=\scriptsize\bfseries,milcTurquesa,anchor=west] at (2.28,0.75)
    {altura $40^\circ$};

  \fill[milcMostaza,draw=milcNegro!55] (Est) circle (3pt);

  % (La síntesis va en el pie de figura vía \guiaDiagrama, para no repetirla
  %  dentro del propio diagrama.)

\end{tikzpicture}
